% PAGES: 92-95
% Continues section 2.8 "Cubic spline interpolation", begun in
% sec02_c6_1.tex (this agent's own file, ends there with eq. (2.79)).

The unknowns are the coefficients $a_i,b_i,c_i,d_i$, $i=1:n$. Thus, there are $4n$
unknowns. There are $n$ constraints corresponding to (2.76), $n-1$ constraints
corresponding to (2.77), $n-1$ constraints corresponding to (2.78), and $n-1$
constraints corresponding to (2.79), for a total of $4n-2$ constraints. Two more
constraints are therefore needed for a unique solution. The most common choice for
these two constraints is to require that $f_1''(x_0)=0$ and $f_n''(x_n)=0$, i.e.,
\begin{align}
2c_1+6d_1x_0 &= 0; \\
2c_n+6d_nx_n &= 0.
\end{align}

The resulting method is called the \textit{natural cubic spline interpolation}.

Let $\bar x$ be the $4n\times1$ vector of the unknowns $a_i,b_i,c_i,d_i$, $i=1:n$,
given by
\[
\bar x(4i-3)=a_i; \quad \bar x(4i-2)=b_i; \quad \bar x(4i-1)=c_i; \quad
\bar x(4i)=d_i; \quad \forall i=1:n.
\]

Note that (2.76--2.81) is a linear system with $4n$ equations and $4n$ unknowns
which can be expressed in matrix notation as
\begin{equation}
\bar M\bar x \;=\; \bar b,
\end{equation}
where $\bar b$ is an $4n\times1$ vector given by
\begin{align}
\bar b(1)&=0; \quad \bar b(4n)=0; \\
\bar b(4i-2)&=v_{i-1}; \quad \bar b(4i-1)=v_i, \quad \forall i=1:n; \\
\bar b(4i)&=0; \quad \bar b(4i+1)=0, \quad \forall i=1:(n-1),
\end{align}
and $\bar M$ is the $4n\times4n$ matrix given by
\begin{align}
\bar M(1,3)&=2; \quad \bar M(1,4)=6x_0; \\
\bar M(4n,4n-1)&=2; \quad \bar M(4n,4n)=6x_n; \\
\bar M(4i-2,4i-3)&=1; \quad \bar M(4i-2,4i-2)=x_{i-1}, \quad \forall i=1:n; \\
\bar M(4i-2,4i-1)&=x_{i-1}^2; \quad \bar M(4i-2,4i)=x_{i-1}^3, \quad \forall i=1:n; \\
\bar M(4i-1,4i-3)&=1; \quad \bar M(4i-1,4i-2)=x_i, \quad \forall i=1:n; \\
\bar M(4i-1,4i-1)&=x_i^2; \quad \bar M(4i-1,4i)=x_i^3, \quad \forall i=1:n; \\
\bar M(4i,4i-2)&=1; \quad \bar M(4i,4i-1)=2x_i, \quad \forall i=1:(n-1); \\
\bar M(4i,4i)&=3x_i^2; \quad \bar M(4i,4i+2)=-1, \quad \forall i=1:(n-1); \\
\bar M(4i+1,4i-1)&=2; \quad \bar M(4i+1,4i)=6x_i, \quad \forall i=1:(n-1); \\
\bar M(4i+1,4i+3)&=-2; \quad \bar M(4i+1,4i+4)=-6x_i, \quad \forall i=1:(n-1).
\end{align}

The matrix $\bar M$ is a banded matrix of band 4. Moreover, the matrix $\bar M$ is
nonsingular for any values of the nodes $x_i$, $i=0:n$.\footnote{While it is
difficult to see that the matrix $\bar M$ is nonsingular by inspecting its
entries, note that the solution to the cubic spline interpolation, which is given
by the solution to the linear system $\bar M\bar x=\bar b$ from (2.82), is
equivalent to finding a solution to the tridiagonal system $Mz=b$ shown in section
6.4. As the matrix $M$ is symmetric positive definite, and therefore nonsingular,
for any values of the nodes $x_i$, $i=0:n$, we conclude that the cubic spline
interpolation has a unique solution for any nodes $x_i$, $i=0:n$. Thus, the linear
system $\bar M\bar x=\bar b$ must have a unique solution, and we can therefore
conclude that the matrix $\bar M$ is nonsingular.}

The linear system $\bar M\bar x=\bar b$ from (2.82) can be solved using the LU with
row pivoting solver from Table 2.11, i.e., $\bar x =
\text{linear\_solve\_lu\_row\_pivoting}(\bar M,\bar b)$.

The cubic spline interpolation pseudocode can be found in Table 2.13.

\begin{table}[H]
\centering
\caption{Pseudocode for the natural cubic spline interpolation}
\fbox{\begin{minipage}{0.85\textwidth}
\textbf{Input:}\\
$x_i = $ interpolation nodes, $i=0:n$\\
$v_i = $ interpolation values, $i=0:n$

\medskip
\textbf{Output:}\\
$a_i,b_i,c_i,d_i = $ cubic polynomials coefficients, $i=1:n$

\medskip
compute vector $\bar b$ from (2.83--2.85)\\
compute matrix $\bar M$ from (2.86--2.95)\\
$\bar x = \text{linear\_solve\_lu\_row\_pivoting}(\bar M,\bar b)$\\
for $i = 1:n$\\
\hspace*{1.5em}$a_i=\bar x(4i-3);\ b_i=\bar x(4i-2);\ c_i=\bar x(4i-1);\
d_i=\bar x(4i);$\\
end
\end{minipage}}
\end{table}

\subsection{Cubic spline interpolation for zero rate curves}

Cubic spline interpolation is often used to obtain a continuous zero rate curve
from values of the zero rate for discrete times which price correctly a set of
interest rate instruments such as bonds and swaps. The zero rate curve thus
obtained can then be used to price interest rate instruments with any cash flow
dates.

For example, assume that the following zero rates have been determined:

\begin{center}
\begin{tabular}{|l|l|}
\hline
Time & Zero Rate \\
\hline
overnight & 0.50\% \\
\hline
2 months & 0.65\% \\
\hline
6 months & 0.85\% \\
\hline
12 months & 1.05\% \\
\hline
20 months & 1.20\% \\
\hline
\end{tabular}
\end{center}

In other words, $r(0,0) = 0.0050$; $r\left(0,\frac2{12}\right) = 0.0065$;
$r\left(0,\frac6{12}\right) = 0.0085$; $r(0,1) = 0.0105$; and
$r\left(0,\frac{20}{12}\right) = 0.0120$, where $r(0,t)$ is the zero rate
corresponding to time $t$.

We use cubic spline interpolation to find the zero rate curve $r(0,t)$ for all
maturities up to 20 months, i.e., for all $0\le t\le\frac{20}{12}$. Thus, $r(0,t)$
is assumed to be a cubic polynomial on each of the intervals
$\left[0,\frac2{12}\right]$, $\left[\frac2{12},\frac6{12}\right]$,
$\left[\frac6{12},1\right]$, and $\left[1,\frac{20}{12}\right]$. The coefficients
of these cubic polynomials are obtained by solving the linear system (2.82)
corresponding to the $16\times16$ banded matrix $\bar M$ given by (2.97--2.98):

\begin{equation}
\bar M \;=\;
{\footnotesize
\left(\begin{array}{cccccccc}
0&0&2&0&0&0&0&0\\
1&0&0&0&0&0&0&0\\
1&\frac16&\frac1{36}&0.0046&0&0&0&0\\
0&1&\frac13&\frac1{12}&0&-1&0&0\\
0&0&2&1&0&0&-2&-1\\
0&0&0&0&1&\frac16&\frac1{36}&0.0046\\
0&0&0&0&1&0.5&\frac14&\frac18\\
0&0&0&0&0&1&1&\frac34\\
0&0&0&0&0&0&2&3\\
0&0&0&0&0&0&0&0\\
0&0&0&0&0&0&0&0\\
0&0&0&0&0&0&0&0\\
0&0&0&0&0&0&0&0\\
0&0&0&0&0&0&0&0\\
0&0&0&0&0&0&0&0\\
0&0&0&0&0&0&0&0
\end{array}\right.}
\quad\cdots
\end{equation}
\begin{equation}
\cdots\quad
{\footnotesize
\left.\begin{array}{cccccccc}
0&0&0&0&0&0&0&0\\
0&0&0&0&0&0&0&0\\
0&0&0&0&0&0&0&0\\
0&0&0&0&0&0&0&0\\
0&0&0&0&0&0&0&0\\
0&0&0&0&0&0&0&0\\
0&0&0&0&0&0&0&0\\
0&-1&0&0&0&0&0&0\\
0&0&-2&-3&0&0&0&0\\
1&0.5&0.25&0.125&0&0&0&0\\
1&1&1&1&0&0&0&0\\
0&1&2&3&0&-1&0&0\\
0&0&2&6&0&0&-2&-6\\
0&0&0&0&1&1&1&1\\
0&0&0&0&1&\frac53&\frac{25}9&4.6296\\
0&0&0&0&0&0&2&10
\end{array}\right)}
\end{equation}

The resulting zero rate curve
\begin{equation}
r(0,t) \;=\;
\begin{cases}
0.005+0.0095t-0.0171t^3, & \text{if } 0\le t\le\frac2{12};\\
0.0049+0.0115t-0.0122t^2+0.0073t^3, & \text{if } \frac2{12}\le t\le\frac6{12};\\
0.0059+0.0057t-0.0006t^2-0.0005t^3, & \text{if } \frac6{12}\le t\le1;\\
0.0044+0.01t-0.0049t^2+0.001t^3, & \text{if } 1\le t\le\frac{20}{12}.
\end{cases}
\end{equation}

The zero rate curve $r(0,t)$ can be used, e.g., to find the value of a 13 months
3\% quarterly bond with face value \$100. This bond has the following cash flows:

\begin{center}
\begin{tabular}{|l|l|}
\hline
Date & Cash flow \\
\hline
1 month & \$0.75 \\
\hline
4 months & \$0.75 \\
\hline
7 months & \$0.75 \\
\hline
10 months & \$0.75 \\
\hline
13 months & \$100.75 \\
\hline
\end{tabular}
\end{center}

The value of the bond, denoted by $B$, is obtained by adding the present value of
all the future cash flows of the bond. Assuming that the zero rate $r(0,t)$
corresponds to continuous compounding, we obtain that
\[
B \;=\; 0.75\exp\left(-\frac1{12}r\left(0,\frac1{12}\right)\right) \;+\;
0.75\exp\left(-\frac4{12}r\left(0,\frac4{12}\right)\right)
\]
\[
+\; 0.75\exp\left(-\frac7{12}r\left(0,\frac7{12}\right)\right) \;+\;
0.75\exp\left(-\frac{10}{12}r\left(0,\frac{10}{12}\right)\right)
\]
\[
+\; 100.75\exp\left(-\frac{13}{12}r\left(0,\frac{13}{12}\right)\right).
\]

From (2.99), we find that $r\left(0,\frac1{12}\right)=0.005782$;
$r\left(0,\frac4{12}\right)=0.007648$; $r\left(0,\frac7{12}\right)=0.008922$;
$r\left(0,\frac{10}{12}\right)=0.009944$; and
$r\left(0,\frac{13}{12}\right)=0.010754$.

Thus, $B=102.5707$, i.e., the value of the bond is \$102.57.

\section{References}

Detailed proofs for the existence of the LU decomposition can be found in Golub
and Van Loan \cite{18}. Trefethen and Bau \cite{43} give an implicit constructive
proof for the existence of the LU decomposition with row pivoting of a
nonsingular matrix, presenting the intuition behind switching the rows of the
lower triangular factor $L$ in the LU with row pivoting algorithm. A discussion of
the stability of the LU decomposition, both from a theoretical and from a
practical standpoint, can be found in Higham \cite{20}.

The numerical solution of the Black--Scholes PDE using finite differences requires
solving linear systems corresponding to the same tridiagonal matrix. For a
classical treatment of the finite difference solution of the Black--Scholes PDE,
see Wilmott \cite{46}; further implementation details can be found in Hirsa
\cite{21}.
